\chapter{总结与展望}
\section{全文总结}
全文的工作分两条线：一条是给不同条件的社区配上合适的控制方法，另一条是做一个能看清每一步决策的管理平台。本节按算法层与平台层两条路径分别整理工作与结论，详细数据与实验证据见第 \ref{chap:5} 章。

\subsection{算法层：双模式协同控制体系}\label{sec:algo-layer}
（一）面向具备专家热力学规则与预测先验的社区场景（对应CHESCA-ResMARL）。分层控制流水线由"负荷预测 → 本地 PID 初稿 → 社区电池 A* 树搜索 → 动作回写"四个环节组成，本文修复了其中生活热水储罐状态更新规划闭环，使储热移峰能力重新生效。在此之上，叠加了一条由当前小时固定负荷脉冲比驱动的连续分级软门控：脉冲越深、对储能充放的允许下调越大。

标准月度（30 天）评估集社区峰值由 12.705 kW 降至 12.119 kW（峰值比由 0.6380 降至 0.6086，相对降幅 4.61\%）。长周期三月（约 92 天）评估集由 12.604 kW 降至 11.260 kW（峰值比降至 0.5256，相对降幅 10.66\%）。在零样本迁移上，长周期三月与六建筑扩容测试集配合用电特征自适应标定分别实现 $-26.18\%$ 与 $-11.11\%$ 的峰值降幅。住户热舒适越界率全程为零。

残差多智能体学习策略在原规则之上叠了三层安全防护：修正幅度限制 $\alpha$、电力通道屏蔽掩码、电池低电量门控。多种子实验证实，四档幅度下热舒适越界变化量均严格保持 $0.000$ 个百分点。放开探索纠偏实验进一步查明：120 轮预算下 $86.2\%$ 的削峰改善来自同策略稳步更新（IPPO）；97 维全局状态拼接的集中式评论家（MAPPO）在 320 轮长训下出现过拟合回退。这一发现与都灵理工大学同类学位论文观察到的现象方向一致 \cite{costea2024}。

（二）面向缺乏专家规则先验的社区场景与工程化落地（对应MARL）。本文设计了一条非侵入式环境包装器，把与建筑数量无关的 9 维区域协同态势特征动态注入各智能体局部观测。奖励侧把"本步向电网买电的电度贡献"分摊惩罚并把目标锚定到运行历史最高峰值；训练侧把经验回放池扩到 200\,000 条，叠目标网络缓慢更新与 3 步多步累积回报引导。

结果上，纯多智能体策略的全时段峰值比由单楼独立 SAC 基线的 0.9196 压至 $0.8872 \pm 0.0172$，热舒适越界率从 5.69\% 降至 3.61\%（相对降幅 36.5\%），并取得所有对比方法中最低的纯电度电费比值（0.8835）。在次高峰时刻（第 27 天晚上 20 时），社区净负荷已低于分层规则基线，说明纯学习在晚高峰前段电池尚有余电时已经能反超。

全月最高用电尖峰时刻（第 28 天下午 22 时）的功率构成分解显示：纯多智能体与分层规则基线之间剩余的 $7.08\text{ kW}$ 差距中，$61.9\%$（$4.384\text{ kW}$）来自峰前日常分时电价套利对电池的消耗，$38.1\%$（$2.693\text{ kW}$）来自稳态空调保底约束限制了大幅关停。这是无预测纯学习策略的物理能力边界。

两条算法线在适用场景与性能边界上互补：CHESCA-ResMARL 适合具备专家规则先验的社区、能在 30 天月度集上把峰值比压到 0.6086，但需要工程化部署分层规则；MARL 适合缺乏规则的社区、稳态峰值比压到 0.8872，但对极端尖峰时刻仍有物理能力上限。

\subsection{平台层：复现护栏、经济核算与可视化管理}\label{sec:plat-layer}
平台层的工作收敛在三件事上：第一件是复现护栏，训练入口的进程哈希种子（\texttt{PYTHONHASHSEED}）自检锁定把同种子跨批次运行的峰值极差从实测的 $0.0894$ 彻底消除为 $0.0000$（详细见第 \ref{sec:reproduce} 节），再用训练末期 3 次检查点均值指标把跨种子标准差再压低约 35\%；第二件是经济核算，建立社区并网点（PCC）两部制最大需量电价与锂电池循环寿命折旧核算模型（式 \eqref{eq:two-part-tariff}），对所有控制策略给出可比的月度电费账单；第三件是可视化管理平台，基于 Vue 前端、Spring Boot 后端、MySQL 数据库与 Python 异步进程池开发 B/S 四层架构，把数据集管理、模型拓扑前置校验、动态分组参数配置、实时训练监控、多维 KPI 对比与逐小时决策过程回放六件事集中到同一界面（系统构建详见第 \ref{chap:4} 章）。其中，协同削峰策略每月仅消耗 6.06 次电池满循环，扣除折旧后仍实现每月净节省 $60.71\text{ \$}$（净降幅 23.9\%），相比纯 SAC 显著下降一个量级且能贴合工商业社区微电网的实际账单结构。

\section{研究不足与未来工作展望}
本章实验留下了三项已经由实测证实的能力边界——它们既是本文的客观短处，也是后续工作的天然起点。第一项边界在全年跨季节评估集上——寒冷冬季住户室内热舒适越界率高达 $49.06\%$，社区峰值改善为零，说明现有控制器仅覆盖空调制冷、热水与电池三通道，还无法应对冬季热泵采暖负荷，必须把采暖通道纳入统一控制维度。第二项边界是纯多智能体策略在建筑规模放大时的零样本性能偏弱：$6$ 栋扩容测试集上，未叠加用电特征自适应标定时仅 $-2.23\%$ 的削峰，必须叠加自适应标定才能稳到 $-11.11\%$，纯多智能体策略在新社区上的迁移代价不容忽视。第三项边界是集中式评论家架构在长训预算下的过拟合风险：表\ref{tab:ctde-ablation} 显示，$320$ 轮训练的 $97$ 维全局评论家峰值回升至 $12.541\text{ kW}$，反不如 $31$ 维轻量同策略 IPPO 的 $12.120\text{ kW}$，高维全局值函数难以稳定匹配社区级策略的价值估计盲区。结合国际同类研究的前沿进展 \cite{javeed2024,costea2024}，后续工作按三个方向推进。

（1）多时间尺度协同学习机制。纯多智能体策略目前只根据历史已发生负荷做反馈调节，平峰与尖峰时段共用同一保底约束——这一短板是引入多时间尺度协同学习的直接动因。本文按三步推进：第一步，将轻量级超短期外生负荷脉冲态势感知特征融入区域协同观测空间；第二步，研究与实时脉冲紧迫度联动的状态自适应安全边界松弛机制——平峰期收紧温控下限保障舒适度，非停电尖峰预警期动态放宽空调柔性减载的上限；第三步，构建"平峰期高电量备电储备—尖峰期归因放电削峰"的多时间尺度奖励解耦方法，目标是突破纯多智能体在极端偶发尖峰时刻的能力上限，摆脱对长周期未来负荷预测器的依赖，同时也缓解建筑规模放大下零样本迁移代价偏高的问题。

（2）控制通道扩展与可微安全层。针对 $49.06\%$ 冬季热舒适越界率与零削峰的边界事实，本文把控制动作维度与安全防护范围扩展到冬季热泵采暖与电动汽车双向充放电（V2G）设备，借鉴隆德大学 Javeed 在展望中提及的受约束马尔可夫决策过程思路 \cite{javeed2024}，把硬裁剪安全边界升级为支持梯度回传的可微安全层。同时引入联邦多智能体强化学习，在不上传住户原始用电曲线的前提下实现跨社区协同训练，应对 $6$ 栋扩容场景下纯 MARL 零样本性能偏弱的问题。

（3）预测辅助在线规划与真实楼宇系统对接。集中式评论家长训的过拟合风险需要从两个方向同时收敛：价值函数设计与轻量评论家逐步扩容。一方面引入更保守的价值正则化，限制全局评论家在高维输入下训练后期的发散空间；另一方面采用"本楼评论家为主、区域摘要为辅"的渐进式架构扩展策略，逐步逼近 $31$ 维同策略 IPPO 的稳定表现。在此基础上把本文的分层负荷预测器与多智能体强化学习的价值网络结合，在实际运行中开展短时前瞻滚动规划，提高策略对极端天气与设备老化的适应能力。管理平台层面，引入 ELK/Loki 集中式日志检索与自动化超参数搜索工具 \cite{costea2024}，通过楼宇工业通信协议（BACnet、Modbus）把平台的安全防护与控制服务对接至真实楼宇自控系统，在实际运行环境中检验算法的长期节能降费效益。

\pagestyle{empty}
